% v1-tate-punct.tex -- M7 evaluation task 3 (G4 tate gate)
% Punctuation rotation behavior: measure physical advance of 、。「」ー
% under CJKvert with the yoko ipaex Type1 font. If vertical writing were
% real, each cell would advance exactly 1zw along the column axis.
% A rotated yoko glyph advances by its HEIGHT (< 1zw for 、。) and the
% glyph lands at the wrong corner of the em box (no vertical variants in
% Type1; rotation is geometric, not glyph-substitution).
\documentclass[paper=a5]{jlreq}
\usepackage{CJKvert}
\makeatletter
\def\probe#1{%
  \setbox0=\hbox{#1}% the glyph as the backend places it (CJKvert rotates inside)
  \typeout{V1-PUNCT \meaning#1: box wd=\the\wd0, ht=\the\ht0, dp=\the\dp0}}
\begin{document}
\probe{あ}
\probe{、}
\probe{。}
\probe{「}
\probe{」}
\probe{ー}
\typeout{V1-PUNCT 1zw = \the\dimexpr1\jlreq@zw\relax}
% kinsoku probe: 、。post-penalties of the backend table are defined for
% the yoko axis only; check they still fire in vertical mode by widening
% a line so a 、falls at margin. (qualitative: compile must rc=0)
あいうえおかきくけこ、さしすせそたちつてと。
「縦書き」のテスト。
\end{document}
